\documentclass[12pt]{article}

% ── Packages ──────────────────────────────────────────────────────────────────
\usepackage[margin=1in]{geometry}
\usepackage{setspace}
\usepackage{booktabs}
\usepackage{tabularx}
\usepackage{array}
\usepackage{amsmath}
\usepackage{amssymb}
\usepackage{hyperref}
\usepackage{xcolor}
\usepackage{titlesec}
\usepackage{enumitem}
\usepackage{parskip}
\usepackage{fancyhdr}
\usepackage{caption}
\usepackage{float}
\usepackage{tikz}
\usetikzlibrary{shapes.geometric, arrows.meta, positioning, fit, backgrounds}

% ── Colors ────────────────────────────────────────────────────────────────────
\definecolor{primaryblue}{RGB}{30, 80, 160}
\definecolor{lightblue}{RGB}{220, 235, 255}
\definecolor{darkgray}{RGB}{60, 60, 60}
\definecolor{midgray}{RGB}{130, 130, 130}

% ── Hyperref ──────────────────────────────────────────────────────────────────
\hypersetup{
    colorlinks=true,
    linkcolor=primaryblue,
    citecolor=primaryblue,
    urlcolor=primaryblue
}

% ── Section formatting ────────────────────────────────────────────────────────
\titleformat{\section}{\large\bfseries\color{primaryblue}}{}{0em}{}[\titlerule]
\titleformat{\subsection}{\normalsize\bfseries\color{darkgray}}{}{0em}{}
\titleformat{\subsubsection}{\normalsize\itshape\color{darkgray}}{}{0em}{}

% ── Header / Footer ───────────────────────────────────────────────────────────
\pagestyle{fancy}
\fancyhf{}
\rhead{\textcolor{midgray}{\small MACS 30113 · Large-Scale Computing · Final Project Proposal}}
\lhead{\textcolor{midgray}{\small Yulin (Olivia) Wang}}
\cfoot{\textcolor{midgray}{\thepage}}
\renewcommand{\headrulewidth}{0.4pt}

% ── Document ──────────────────────────────────────────────────────────────────
\begin{document}

\begin{center}
    {\LARGE\bfseries\color{primaryblue}
        Do LLMs Read 10-Ks Better Than Dictionaries?}\\[6pt]
    {\large A Staggered Event Study of Financial Text Measures\\
    and Earnings Announcement Returns}\\[18pt]
    {\normalsize Yulin (Olivia) Wang}\\
    {\small MAPSS, University of Chicago}\\[4pt]
    {\small\textcolor{midgray}{MACS 30113: Large-Scale Computing in the Social Sciences}}\\
    {\small\textcolor{midgray}{Final Project Proposal · \today}}
\end{center}

\vspace{0.5em}
\hrule height 1.5pt
\vspace{1.5em}

% ── 1. Motivation ─────────────────────────────────────────────────────────────
\section{Motivation and Research Question}

Annual filings (10-Ks) submitted to the SEC contain rich qualitative information
about firm performance, risk outlook, and managerial tone. Since
\citet{LoughranMcDonald2011}, the standard approach for extracting sentiment from
financial text has been dictionary-based counting: each document is scored by the
fraction of words belonging to pre-defined word lists (negative, uncertain,
litigious, etc.). While tractable at scale, this method is bag-of-words by
construction and cannot capture context, negation, or domain-specific nuance.

Large language models (LLMs) offer a fundamentally different approach: rather
than counting keywords, they process entire passages and produce semantic
judgments about tone, uncertainty, and forward-looking orientation. A growing
body of work suggests LLMs outperform dictionaries in short-window sentiment
classification tasks \citep{LopezLira2023}, but systematic evidence using a
rigorous causal identification framework at scale remains limited.

\medskip
\noindent\textbf{Research Question:} Does LLM-based sentiment scoring of 10-K
filings predict earnings-announcement abnormal returns more precisely than the
Loughran--McDonald (L\&M) dictionary method, and does this advantage vary
systematically across firm size, industry, and macroeconomic environment?

% ── 2. Contribution ───────────────────────────────────────────────────────────
\section{Contribution}

This project makes three contributions relative to the existing literature:

\begin{enumerate}[leftmargin=1.5em, label=\textbf{\arabic*.}]
    \item \textbf{Identification strategy.} Existing comparisons of LLM and
    dictionary measures rely on simple OLS regressions. This paper uses a
    \emph{staggered event study} design, exploiting the staggered timing of
    10-K filing dates across firms and years to estimate causal effects on
    cumulative abnormal returns (CARs) in a difference-in-differences spirit.

    \item \textbf{Scale.} The full analysis covers approximately 50,000--80,000
    10-K filings (2010--2020) across S\&P 1500 constituents, processed via GPU-
    accelerated batch inference on a high-performance computing cluster. This
    scale is necessary to achieve sufficient statistical power for heterogeneity
    analysis.

    \item \textbf{Multi-dimensional comparison.} Rather than a single sentiment
    score, both methods are compared along three parallel dimensions: overall
    \emph{tone}, \emph{uncertainty}, and \emph{forward-looking intensity},
    enabling a granular assessment of where each approach gains or loses
    informational content.
\end{enumerate}

% ── 3. Data ───────────────────────────────────────────────────────────────────
\section{Data}

\begin{table}[H]
\centering
\caption{Data Sources and Variables}
\renewcommand{\arraystretch}{1.35}
\begin{tabularx}{\textwidth}{l l X}
\toprule
\textbf{Variable} & \textbf{Source} & \textbf{Details} \\
\midrule
10-K full text & SEC EDGAR & Full-text search API
    (\texttt{efts.sec.gov}); MD\&A section extracted via regex \\
Filing dates \& CIK--ticker map & WRDS / Compustat & \texttt{date\_filed},
    matched to CRSP permno via CRSP--Compustat link table \\
Daily stock returns & WRDS / CRSP & \texttt{ret}, value-weighted market
    return \texttt{vwretd} for excess-return calculation \\
Firm fundamentals & WRDS / Compustat & Total assets (\texttt{at}), SIC code,
    book-to-market ratio \\
\bottomrule
\end{tabularx}
\end{table}

\noindent\textbf{Sample.} S\&P 1500 constituent firms, fiscal years 2010--2020.
Each observation is a firm-year 10-K filing linked to the corresponding
earnings-announcement trading window. Firms with fewer than 3 years of
consecutive data are dropped.

% ── 4. Methodology ────────────────────────────────────────────────────────────
\section{Methodology}

The analysis proceeds in two sequential steps.

\subsection{Step 1: Text Measurement (Machine Learning)}

\subsubsection{Track A — L\&M Dictionary (Baseline Replicate)}

Following \citet{LoughranMcDonald2011}, each 10-K is scored as:
\begin{equation}
    \text{Negative Ratio}_{it} = \frac{\text{\# negative words}}{\text{total words}}, \quad
    \text{Uncertainty Ratio}_{it} = \frac{\text{\# uncertain words}}{\text{total words}}
\end{equation}
Three measures are constructed per filing: \texttt{negative\_ratio},
\texttt{uncertainty\_ratio}, and \texttt{litigious\_ratio}. These require only
CPU computation and serve as the benchmark.

\subsubsection{Track B — LLM Scoring (Extension)}

The MD\&A section of each 10-K (typically 2,000--5,000 words) is passed to an
instruction-tuned LLM for structured inference. The candidate model is
\textbf{Llama-3-8B-Instruct} (or FinBERT as a lighter alternative), deployed
via GPU batch inference on the Midway3 HPC cluster. The inference prompt
requests a JSON-formatted output with three scores:

\begin{center}
\texttt{\{"tone": [0,1], "uncertainty": [0,1], "forward\_looking": [0,1]\}}
\end{center}

where each dimension is scored on a continuous $[0,1]$ scale, operationalized
as follows:

\begin{itemize}[leftmargin=1.5em, noitemsep]
    \item \textbf{Tone}: overall positivity of managerial language ($0=$ strongly
    negative, $1=$ strongly positive)
    \item \textbf{Uncertainty}: degree of hedging and epistemic qualification
    \item \textbf{Forward-looking Intensity}: density of prospective statements
    about future performance
\end{itemize}

Batch inference is parallelized across SLURM job arrays (5,000 filings per
array element), with AWS free-tier GPU instances used as overflow capacity.
Expected throughput: $\approx$10,000 documents per GPU-hour.

\subsection{Step 2: Staggered Event Study (Empirical Identification)}

\subsubsection{Abnormal Return Calculation}

For each filing event $(i, t)$, cumulative abnormal returns (CARs) are computed
over windows $[-1,+1]$ and $[-3,+3]$ relative to the filing date:
\begin{equation}
    \text{CAR}_{it}[\tau_1, \tau_2] = \sum_{\tau=\tau_1}^{\tau_2}
    \bigl(r_{it\tau} - \hat{r}_{it\tau}\bigr)
\end{equation}
where $\hat{r}_{it\tau} = \hat{\alpha}_i + \hat{\beta}_i r_{mt\tau}$ is the
market-model expected return estimated over the pre-event window $[-250, -11]$.

\subsubsection{Main Regression}

The primary specification is:
\begin{equation}
    \text{CAR}_{it} = \beta_1 \cdot \text{LLM\_tone}_{it}
                    + \beta_2 \cdot \text{L\&M\_negative}_{it}
                    + \gamma' \mathbf{X}_{it}
                    + \alpha_i + \delta_t + \varepsilon_{it}
    \label{eq:main}
\end{equation}
where $\alpha_i$ are firm fixed effects, $\delta_t$ are year-quarter fixed
effects, $\mathbf{X}_{it}$ is a vector of controls (log assets, book-to-market,
leverage), and standard errors are clustered at the firm level.

\subsubsection{Horse-Race Specification}

To directly test the marginal predictive content of LLM scoring over the
dictionary approach, both measures enter equation~\eqref{eq:main} simultaneously.
A larger and more precisely estimated $\hat{\beta}_1$ relative to $\hat{\beta}_2$
is interpreted as evidence that LLM scoring provides incremental information
beyond word counting.

% ── 5. Heterogeneity ──────────────────────────────────────────────────────────
\section{Heterogeneity Analysis}

Three pre-registered dimensions of heterogeneity are examined by interacting the
text measures with indicator variables:

\begin{enumerate}[leftmargin=1.5em, noitemsep, label=\textbf{H\arabic*.}]
    \item \textbf{Firm size.} Small-cap vs.\ large-cap firms (split at NYSE
    median market capitalization). LLM measures are hypothesized to show larger
    gains for small firms, where analyst coverage is sparse and qualitative
    disclosure carries more incremental weight.

    \item \textbf{Industry.} High-complexity sectors (technology, biotech,
    finance) vs.\ low-complexity sectors (utilities, materials). Semantic
    complexity may amplify the dictionary's limitations and LLM's advantages.

    \item \textbf{Macroeconomic environment.} Crisis periods (2010--2011
    recovery, 2020 COVID shock) vs.\ stable periods. Uncertainty language may
    be particularly informative during high-volatility regimes.
\end{enumerate}

% ── 6. Computational Architecture ─────────────────────────────────────────────
\section{Computational Architecture and Timeline}

\begin{table}[H]
\centering
\caption{Four-Week Project Timeline}
\renewcommand{\arraystretch}{1.35}
\begin{tabularx}{\textwidth}{c l X}
\toprule
\textbf{Week} & \textbf{Focus} & \textbf{Tasks} \\
\midrule
1 & Data acquisition
  & SEC EDGAR bulk download (10-K full texts); WRDS pull for filing dates,
    CRSP returns, and Compustat fundamentals; data cleaning and CIK--permno
    merge \\
2 & Text measurement
  & L\&M dictionary scoring (CPU); LLM batch inference via Midway3 SLURM job
    arrays; quality checks on JSON outputs; score distribution validation \\
3 & Empirical analysis
  & Market-model CAR estimation; main regression (Eq.\ \ref{eq:main});
    horse-race specification; robustness checks (alternative windows,
    winsorization) \\
4 & Heterogeneity \& write-up
  & Subgroup regressions (H1--H3); visualization; final report and
    presentation slides \\
\bottomrule
\end{tabularx}
\end{table}

\noindent\textbf{Computational resources:}
\begin{itemize}[leftmargin=1.5em, noitemsep]
    \item \textbf{Midway3 HPC} (University of Chicago): Primary GPU inference
    via SLURM job arrays. Target allocation: A100 or V100 nodes, $\sim$20
    GPU-hours total.
    \item \textbf{AWS (free-tier, \$50 credit)}: Overflow parallel compute for
    robustness checks and alternative LLM variants.
    \item \textbf{Local / Notebook}: Data wrangling, L\&M scoring, regression
    analysis (\texttt{pyfixest}, \texttt{linearmodels}).
\end{itemize}

% ── 7. Expected Results ───────────────────────────────────────────────────────
\section{Expected Results and Significance}

Based on prior literature, the following outcomes are anticipated:

\begin{itemize}[leftmargin=1.5em, noitemsep]
    \item \textbf{Replicate:} L\&M negative ratio significantly predicts
    $\text{CAR}[-1,+1]$, consistent with \citet{LoughranMcDonald2011}.
    \item \textbf{Extension:} LLM-based tone exhibits a larger and more precisely
    estimated coefficient in the horse-race specification, particularly for
    short-window CARs where information release is concentrated.
    \item \textbf{Heterogeneity:} LLM advantage is concentrated among small-cap
    firms and high-complexity industries, consistent with the hypothesis that
    semantic depth matters most where dictionary noise is highest.
\end{itemize}

These findings would contribute actionable guidance for practitioners and
researchers choosing between text-analysis paradigms, and methodologically
demonstrate the value of combining large-scale LLM inference with rigorous
econometric identification.

% ── 8. References ─────────────────────────────────────────────────────────────
\section*{References}
\addcontentsline{toc}{section}{References}

\begin{thebibliography}{9}

\bibitem[Loughran and McDonald(2011)]{LoughranMcDonald2011}
Loughran, T., \& McDonald, B. (2011).
\newblock When is a liability not a liability? Textual analysis, dictionaries,
and 10-Ks.
\newblock \textit{Journal of Finance}, 66(1), 35--65.

\bibitem[Lopez-Lira et al.(2023)]{LopezLira2023}
Lopez-Lira, A., Tang, Y., \& Ye, Z. (2023).
\newblock Can ChatGPT forecast stock price movements?
\newblock \textit{SSRN Working Paper}.

\bibitem[Ke et al.(2019)]{KeKellyXiu2019}
Ke, Z.~T., Kelly, B.~T., \& Xiu, D. (2019).
\newblock Predicting returns with text data.
\newblock \textit{NBER Working Paper No. 26186}.

\bibitem[Huang et al.(2018)]{Huang2018}
Huang, A.~H., Lehavy, R., Zang, A.~Y., \& Zheng, R. (2018).
\newblock Analyst information discovery and interpretation roles: A topic
modeling approach.
\newblock \textit{Management Science}, 64(6), 2833--2855.

\end{thebibliography}

\end{document}
